\booksection{Lampiran: Analisis ekspansi memori ponsel Android}

\chapterintro{Saat membeli ponsel, Anda sering melihat promosi seperti “12GB+8GB RAM”—seperti mendapat tambahan 8GB RAM secara gratis. Tapi apakah benar bisa membuat ponsel lebih cepat? Apakah justru merusak ponsel? Di bagian ini kita akan jelaskan teknologi “ekspansi memori” secara menyeluruh.}

\booksubsection{Definisi esensi}

“Ekspansi memori” ponsel Android (juga disebut memory extension, memory fusion, RAM Plus, dll)\textbf{bukan benar-benar menambah chip RAM secara fisik}, melainkan teknologi software system-level: dengan mengalokasikan secara dinamis sebagian ruang penyimpanan internal ponsel (ROM/UFS flash) untuk disimulasikan sebagai virtual RAM.

Sebagai perumpamaan: RAM fisik adalah “meja kerja kecepatan tinggi”, ekspansi memori adalah “gudang penyimpanan sementara”. Saat meja kerja tidak cukup, sistem akan memindahkan data yang sementara tidak digunakan ke gudang, dan dipindahkan kembali saat dibutuhkan. Gudang kapasitasnya besar, tapi bolak-balik pasti lebih lambat daripada langsung di meja kerja.

\booksubsection{Prinsip teknologi}

Implementasi produsen mainstream utamanya ada dua cara:

\subsubsection*{Kompresi zRAM}

Membagi ruang data terkompresi di dalam RAM, menyimpan kode aplikasi tidak aktif secara terkompresi, rasio kompresi sekitar 2:1. Data masih di RAM, kecepatan cepat, tidak mengurangi flash, tapi butuh CPU untuk kompresi-dekompresi.

\subsubsection*{Partisi swap}

Langsung membagi ruang tetap di UFS flash sebagai virtual memory, menulis data proses latar belakang ke penyimpanan. Bergantung pada bandwidth penyimpanan, kecepatan baca-tulis jauh lebih rendah dari RAM fisik.

Promosi produsen “12GB+8GB RAM” biasanya adalah kombinasi 12GB RAM fisik + 8GB virtual memory flash. Perlu dicatat, zRAM default tidak menempati penyimpanan bodi, hanya saat zRAM sepenuhnya terisi data, sistem baru akan membuat swap file di penyimpanan, dan situasi ini sangat jarang terjadi.

\booksubsection{Kesenjangan kecepatan adalah bottleneck kunci}

Keterbatasan mendasar teknologi ini adalah perbedaan kecepatan. UFS 4.0 flash mainstream saat ini sequential read speed sekitar 4200MB/s, hanya 1/10–1/15 dari RAM fisik LPDDR5X. Karena itu pada tahap kompresi zRAM beban ringan kecepatan mendekati RAM nyata, tapi begitu masuk skenario beban tinggi, kecepatan akan turun drastis ke puluhan kali lebih lambat.

\booksubsection{Skenario yang cocok dan batas manfaat}

Ekspansi memori bukan panacea, efek berbeda tergantung model dan skenario penggunaan:

\begin{itemize}
  \item \textbf{Skenario bermanfaat}: Model dengan 6GB RAM fisik + UFS 3.1 ke atas, ruang tersisa $\geq$30GB, setelah mengaktifkan ekspansi 3GB jumlah keep-alive latar belakang bisa lebih 3-4, reload rate turun 68\%, cocok untuk pengguna multi-tasking berat seperti kurir delivery, driver ride-hailing.
  \item \textbf{Skenario harian}: Peningkatan signifikan untuk multi-akun WeChat, keep-alive tab browser, respons pesan instant messaging, waktu cold start aplikasi latar belakang jelas lebih singkat.
  \item \textbf{Skenario terbatas}: Peningkatan terbatas untuk skenario bandwidth tinggi seperti loading game besar, editing video 4K.
  \item \textbf{Skenario merugikan}: Model dengan 8GB RAM + UFS 2.2 dan ruang tersisa tidak cukup, setelah diaktifkan fluktuasi frame rate game justru membesar, scroll video pendek lebih boros daya, switch lag.
  \item \textbf{Skenario tidak efektif}: Model RAM besar di atas 16GB manfaatnya hampir nol, paling banyak keep-alive 1 aplikasi lagi, tapi akan membuat beban CPU naik 12\%, suhu naik 2-3℃. Untuk model RAM di atas 12GB, disarankan langsung dimatikan untuk menghindari komputasi kompresi yang tidak bermakna.
\end{itemize}

\booksubsection{Penamaan dan cara mengaktifkan tiap produsen}

Penamaan berbeda-beda antar merek: Xiaomi menyebut “Memory Extension”, Huawei “Memory Fusion”, Samsung “RAM Plus”, vivo “Extended RAM”, OnePlus “Memory Extension”. Fitur ini adalah kemampuan native sistem, tidak terikat langsung dengan versi underlying Android.

Path umum mengaktifkan: Masuk [Pengaturan] [Pengaturan lainnya] temukan submenu [Memori] atau [Performa] aktifkan “Ekspansi memori” dan pilih tier kapasitas (biasanya 2GB–7GB), berlaku setelah reboot. Setelah diaktifkan sistem akan menampilkan kombinasi memori seperti “8GB+4GB”.

\begin{tcbwarn}{Pengaturan dalam mode pengembang dapat mempengaruhi stabilitas sistem, opsi yang tidak dipahami jangan diubah sembarangan. Setelah mengaktifkan debugging USB, berhati-hatilah saat menghubungkan ke komputer, jangan hubungkan ke komputer yang tidak dipercaya.}
\begin{itemize}
  \item Ekspansi memori akan menempati sebagian ruang penyimpanan bodi, pastikan penyimpanan tersisa tidak kurang dari 20-30GB, jika tidak bisa mempengaruhi kelancaran sistem.
  \item Flash sebagai virtual memory akan ada keausan baca-tulis, meskipun UFS flash modern umur cukup panjang, tapi swap frekuensi tinggi jangka panjang tetap bisa mempercepat penuaan penyimpanan.
  \item Jika RAM fisik ponsel kecil (seperti 4GB ke bawah), mengaktifkan ekspansi memori hanya bisa meringankan masalah sesaat, tidak bisa menyelesaikan masalah lag secara mendasar.
\end{itemize}

\end{tcbwarn}

\begin{tcbinfo}{Rangkuman}
Ekspansi memori adalah teknologi optimisasi system-level yang matang, bisa meningkatkan kemampuan multi-tasking dan rate keep-alive latar belakang model mid-low end tanpa mengubah hardware, tapi terbatas oleh kesenjangan kecepatan besar antara flash dan RAM fisik, bukan “obat ajaib universal”. Apakah diaktifkan, seberapa besar, perlu dipertimbangkan berdasarkan ukuran RAM fisik model, spesifikasi flash, ruang penyimpanan tersisa dan kebiasaan penggunaan.

\end{tcbinfo}

\begin{tcbpractice}{Mari praktik: Lihat apakah ponsel Anda punya ekspansi memori}
\begin{enumerate}
  \item Buka [Pengaturan], cari "Memory Extension", "Memory Fusion", "RAM Plus" atau "Extended RAM".
  \item Setelah menemukan, lihat apakah saat ini aktif, dan berapa pengaturan kapasitas ekspansi.
  \item Sesuaikan dengan konten bagian ini untuk menilai: Berapa RAM fisik model Anda? Spesifikasi flash apa? Apakah ruang penyimpanan tersisa cukup?
  \item Coba matikan atau turunkan kapasitas ekspansi, amati perubahan keep-alive latar belakang dan panas selama satu minggu, temukan pengaturan yang paling cocok.
\end{enumerate}

\end{tcbpractice}
